\section{Graphene tight-binding \texorpdfstring{$\rightarrow$}{->} Dirac in \texorpdfstring{$2{+}1$}{(2+1)}}
\label{p:graphene}

Graphene~\cite{Semenoff1984,CastroNeto2009} gives the cleanest
two-dimensional version. Working with nearest-neighbor vectors
\[
\vdelta_1 = (a/2)(1,\sqrt{3}),\qquad
\vdelta_2 = (a/2)(1,-\sqrt{3}),\qquad
\vdelta_3 = -a\,(1,0) ,
\]
the structural factor is
\begin{equation}
\Phi(\vk) \;=\;
\sum_{i=1}^{3} \mathrm{e}^{\ii\,\vk\cdot\vdelta_{i}}
\;=\; \mathrm{e}^{\ii\,k_x a/2}\;2\cos(\sqrt{3}\,k_y a / 2)
  \;+\; \mathrm{e}^{-\ii k_x a},
\label{p:eq:Phi}
\end{equation}
and the bipartite Hamiltonian in the $(A,B)$ basis is
\begin{equation}
\Hop_{\mathrm{TB}}(\vk) \;=\; -t\!\begin{pmatrix} 0 & \Phi(\vk) \\ \Phi^{*}(\vk) & 0 \end{pmatrix}.
\end{equation}

\paragraph{Dirac points and effective cone.}
$\Phi$ vanishes at the corners of the Brillouin zone, one of which is
$\vk_{K}=\bigl(2\pi/(3a),\,2\pi/(3\sqrt{3}a)\bigr)$. Verified directly:
$\Phi(\vk_{K})=0$
(\texttt{validators/graphene.py::test\_phi\_vanishes\_at\_K}).
The leading non-trivial behavior is captured by the gradient of $\Phi$ at
$\vk_{K}$, whose two components $c_x = (\partial_{k_x}\Phi)(\vk_{K})$ and
$c_y = (\partial_{k_y}\Phi)(\vk_{K})$ satisfy
\begin{equation}
\lvert c_x\rvert^{2}\;=\;\lvert c_y\rvert^{2}\;=\;\bigl(3a/2\bigr)^{2},
\qquad
\operatorname{Re}\!\bigl(c_x^{*}c_y\bigr)\;=\;0 .
\label{p:eq:c-magnitudes}
\end{equation}
The first identity fixes the magnitude of the Fermi velocity; the second
guarantees that the bilinear $\lvert\Phi(\vk_{K}+\vq)\rvert^{2}$ is
diagonal in $(q_x,q_y)$. Together they give
\begin{equation}
\lvert\Phi(\vk_{K}+\vq)\rvert^{2}\;=\;\tfrac{9}{4}\,a^{2}\,(q_x^{2}+q_y^{2})
+ O(\vq^{3}),
\end{equation}
i.e.\ an isotropic Dirac cone with $v_{F}=\tfrac{3}{2}\,t\,a$.

\paragraph{A computational note.}
The naive verification --- expand $\lvert\Phi\rvert^{2}$ symbolically and
match coefficients --- is computationally expensive because the bilinear
involves nine cross terms with complex exponentials and two-variable
series expansions. The gradient-based verification of
\eqref{p:eq:c-magnitudes} computes only two derivatives and finishes in
milliseconds. We mention this because the analogous reformulation may
help in any symbolic verification of band-crossing structures in higher
dimensions or on more complex Bravais lattices.
\begin{description}
\item[Verified:]
\texttt{validators/graphene.py::test\_gradient\_modulus\_is\_dirac},
\texttt{::test\_cross\_term\_is\_imaginary},
\texttt{::test\_fermi\_velocity}.
\end{description}
